% Part: lambda-calculus
% Chapter: introduction
% Section: unique-readability

\documentclass[../../../include/open-logic-section]{subfiles}

\begin{document}

\olfileid{lam}{syn}{unq}
\olsection{ఏకార్థ పఠనీయత}

ప్రతి పదాన్ని నిర్మించడానికి ఒకే మార్గం ఉందా అని అడగవచ్చు; నిజంగానే
ఉంది. ప్రతి లాంబ్డా పదాన్ని నిర్మించడానికీ, దాన్ని అర్థం చేసుకోవడానికీ
ఒక్కటే మార్గం ఉంది. కింది చర్చలో, \emph{నిర్మాణం} అంటే
\olref[trm]{defn:term}లోని నిర్మాణ నియమాలను ఒకసారి లేదా అనేకసార్లు
ఉపయోగించి పదాన్ని ఏర్పరచే ప్రక్రియ.

\begin{lem}\ollabel{lem:term-start}
పదం ఒక చరంతో లేదా ఎడమ కుండలీకరణ చిహ్నంతో మొదలవుతుంది.
\end{lem}

\begin{proof}
\olref[trm]{defn:term} ప్రకారం నిర్మించినదే పదం అవుతుంది.
\olref[trm]{defn:term-var} ఫలితమైతే అది చరం అయి ఉండాలి.
\olref[trm]{defn:term-abs} లేదా \olref[trm]{defn:term-app} ఫలితమైతే,
ఎడమ కుండలీకరణ చిహ్నంతో మొదలవుతుంది.
\end{proof}

\begin{lem}\ollabel{lem:app-start}
ప్రయోగం ఫలితం రెండు ఎడమ కుండలీకరణ చిహ్నాలతో లేదా ఒక ఎడమ
కుండలీకరణ చిహ్నం తరువాత ఒక చరంతో మొదలవుతుంది.
\end{lem}

\begin{proof}
$M$ ప్రయోగం ఫలితమైతే, అది $(PQ)$ రూపంలో ఉంటుంది; కాబట్టి ఎడమ
కుండలీకరణ చిహ్నంతో మొదలవుతుంది. $P$ ఒక పదం కనుక
\olref{lem:term-start} ప్రకారం అది ఎడమ కుండలీకరణ చిహ్నంతో లేదా
చరంతో మొదలవుతుంది.
\end{proof}

\begin{lem}\ollabel{lem:initial}
ఒక పదం యొక్క ఏ నిజ ఆద్యఖండమూ స్వయంగా పదం కాదు.
\end{lem}

\begin{prob}
పదాల పొడవుపై ఆగమనం ద్వారా
\olref[lam][syn][unq]{lem:initial}ను నిరూపించండి.
\end{prob}

\begin{prop}[ఏకార్థ పఠనీయత] \ollabel{prop:unq}
ప్రతి పదానికి ఒక్కటే నిర్మాణం ఉంటుంది. మరో మాటలో, ఒక నిర్మాణం ద్వారా
పదం~$M$ ఏర్పడితే, అదే పదాన్ని ఏర్పరచగల నిర్మాణం మరొకటి ఉండదు.
\end{prop}

\begin{proof}
  పదాల నిర్మాణంపై ఆగమనం ద్వారా నిరూపిస్తాం.
  \begin{enumerate}
    \item $M$ అనేది $x$ రూపంలో ఉండనివ్వండి; ఇక్కడ $x$ ఒక చరం.
      అమూర్తీకరణ, ప్రయోగం ఫలితాలు ఎప్పుడూ ఎడమ కుండలీకరణ చిహ్నంతో
      మొదలవుతాయి కనుక అవి $M$ను నిర్మించి ఉండలేవు. అందువల్ల $M$
      నిర్మాణం \olref[trm]{defn:term}\olref[trm]{defn:term-var}లోని
      ఒక్క దశ మాత్రమే.
    \item $M$ అనేది $(\lambd[x][N])$ రూపంలో ఉండనివ్వండి; ఇక్కడ $x$
      ఒక చరం, $N$ ఒక పదం. ఇది ఒక్క చరం కాదు కనుక
      \olref[trm]{defn:term}\olref[trm]{defn:term-var} ప్రకారం
      ఏర్పడి ఉండదు. \olref{lem:app-start} ప్రకారం ఇది ప్రయోగం
      ఫలితమూ కాదు. కాబట్టి $M$, $N$పై అమూర్తీకరణ ద్వారా మాత్రమే
      వస్తుంది. ఆగమన పరికల్పన ప్రకారం $N$ నిర్మాణం కూడా ఏకైకం.
    \item $M$ అనేది $(PQ)$ రూపంలో ఉండనివ్వండి; ఇక్కడ $P$, $Q$
      పదాలు. ఇది ఎడమ కుండలీకరణ చిహ్నంతో మొదలవుతుంది కనుక
      \olref[trm]{defn:term}\olref[trm]{defn:term-var} ప్రకారం
      కూడా ఏర్పడి ఉండదు. \olref{lem:term-start} ప్రకారం $P$,
      $\lambd$తో మొదలుకాదు; అందువల్ల $(PQ)$ అమూర్తీకరణ ఫలితం
      కాలేదు. ఇప్పుడు ప్రయోగం ద్వారానే $M$ను నిర్మించే మరో మార్గం
      ఉందని, ఉదాహరణకు అది $(P'Q')$ రూపంలో కూడా ఉందని అనుకుందాం.
      అప్పుడు $P$, $P'$కు నిజ ఆద్యఖండం, లేదా దీనికి విరుద్ధంగా
      ఉండాలి; \olref{lem:initial} ప్రకారం
      ఇది అసాధ్యం. కనుక $P$, $Q$ ఏకైకంగా నిర్ణయితమవుతాయి;
      ఆగమన పరికల్పన ప్రకారం $P$, $Q$ల నిర్మాణాలూ ఏకైకమే.
  \end{enumerate}
\end{proof}

పైనున్న ప్రతిపాదనను మరింత సులభంగా ఇలా చెప్పవచ్చు:
\begin{prop}
  పదం~$M$ కింది రూపాల్లో ఒకదానిలో మాత్రమే ఉండగలదు:
  \begin{enumerate}
    \item $x$; ఇక్కడ $x$ అనేది $M$ ద్వారా ఏకైకంగా నిర్ణయితమయ్యే చరం.
    \item $(\lambd[x][N])$; ఇక్కడ $x$ ఒక చరం, $N$ మరొక పదం;
      రెండూ $M$ ద్వారా ఏకైకంగా నిర్ణయితమవుతాయి.
    \item $(PQ)$; ఇక్కడ $P$, $Q$ అనే రెండు పదాలూ $M$ ద్వారా
      ఏకైకంగా నిర్ణయితమవుతాయి.
  \end{enumerate}
\end{prop}

\end{document}
